\documentclass[10pt,a4paper]{amsart}
\usepackage[margin=19mm]{geometry}
\usepackage{amsmath,amssymb,mathtools}
\usepackage[scr=rsfs,cal=euler]{mathalfa}
\usepackage{xfrac}
\usepackage[unicode,hidelinks,bookmarks=false]{hyperref}
\newcommand{\ie}{i.e.\ }
\newcommand{\cf}{cf.\ }
\newcommand{\resp}{respectively\ }
\newcommand{\Subsection}{\subsection}
\providecommand{\nobreakdash}{\nobreak\hbox{-}}
\setlength{\emergencystretch}{2em}
\multlinegap0pt
\usepackage{amssymb}
\newtheorem{lem}{Lemma}
\DeclareMathOperator{\Id}{Id}
\newcommand{\kx}{k^\times}
\title{Skew fields not finitely generated as algebras}
\author{K. R. Goodearl, E. S. Letzter}
\date{Source: arXiv:2604.18480v2 (CC BY 4.0)}
\begin{document}
\maketitle

\noindent\emph{Source and licence.} Skew fields not finitely generated as algebras. Version arXiv:2604.18480v2 (CC BY 4.0), distributed by arXiv under the Creative Commons Attribution 4.0 International licence, http://creativecommons.org/licenses/by/4.0/, as recorded in the arXiv licence metadata for this exact version and shown on the abstract page https://arxiv.org/abs/2604.18480v2. Full licence notice: \texttt{licenses/arxiv-2604.18480-CC-BY-4.0.txt}.\par\smallskip

\noindent\textit{Editorial scope.} Counted unit: the article's lem iso.simples.vs.eigenvals (source0.tex lines 266--274). The article's complete original proof of it is delivered with it (source0.tex lines 276--284, one \texttt{proof} environment of 9 lines). No mathematics is written, repaired or re-derived: the statement and the proof are the article's own lines, in the article's own notation, and no retained line is substituted or re-flowed.  No further source-local context is needed: every cross-reference of every retained line resolves inside the two delivered spans.  Nothing was omitted from any retained span, so the package carries no omission record.  No retained line contains a citation, so the wrapper carries no bibliography and no bibliography entry.  No retained line contains a figure environment, an image-inclusion command or drawing code, so no image is delivered and the package declares no asset dependency.  Editorial formatting only, in addition: an AMS \texttt{amsart} preamble that restates the article's own declarations and macros used by the retained lines, namely the environments \texttt{lem} on separate counters, together with the standard packages those lines need.  The retained environments print in this document as Lemma 1 in order of appearance, whereas the article prints the same items with the numbering of its own full document.  The complete native source of the article as distributed in the arXiv e-print for this exact version is bundled unchanged as \texttt{sources/198803/source0.tex}.
\begin{lem}  \label{iso.simples.vs.eigenvals}
Let $R = D[y;\tau,\delta]$ where $D$ is a division algebra over a field $k$.

{\rm(a)} Let $a,b \in D$.  Then $R/(y-a)R \cong R/(y-b)R$ if and only if there is some nonzero $v\in D$ such that $av-\tau(v)b = \delta(v)$.

{\rm(b)} Assume that $\delta = 0$, and let $\alpha,\beta \in \kx$.  Then $R/(y-\alpha)R \cong R/(y-\beta)R$ if and only if $\alpha\beta^{-1}$ is an eigenvalue for $\tau$.

{\rm(c)} Assume that $\tau = \Id_D$, and let $\alpha,\beta \in k$. Then $R/(y-\alpha)R \cong R/(y-\beta)R$ if and only if $\alpha - \beta$ is an eigenvalue for $\delta$.
\end{lem}

\begin{proof}
For any $c\in D$, the right $R$-module $R/(y-c)R$ is simple, being $1$-di\-men\-sional over $D$ on the right.

(a) If $R/(y-a)R \cong R/(y-b)R$, there is a nonzero element $x \in R/(y-a)R$ such that $x(y-b) = 0$.  We may write $x = u+(y-a)R$ for some nonzero $u \in D$, and then $u(y-b) = (y-a)v$ for some nonzero $v \in R$.  Since $u(y-b)$ has degree $1$, $v$ lies in $D$.  Comparing coefficients, we find that $u = \tau(v)$ and $- ub = \delta(v) - av$, whence $av-\tau(v)b = \delta(v)$.  The converse is obtained by reversing the above steps.

(b) In view of part (a), the isomorphism occurs if and only if there is some nonzero $v \in D$ such that $\alpha v = \tau(v) \beta$, that is, $\tau(v) = \alpha \beta^{-1} v$.

(c) In this case, the isomorphism occurs if and only if there is some nonzero $v \in D$ such that $\delta(v) = (\alpha-\beta)v$.
\end{proof}

\end{document}
